\documentclass[lettersize,journal]{IEEEtran}
\usepackage{amsmath,amsfonts}
\usepackage{algorithmic}
\usepackage{algorithm}
\usepackage{array}
\usepackage[caption=false,font=footnotesize]{subfig}
\usepackage{textcomp}
\usepackage{stfloats}
\usepackage{url}
\usepackage{verbatim}
\usepackage{graphicx}
\usepackage{cite}
\usepackage{amsmath}
\usepackage{amssymb}
\usepackage{amsthm}

\newtheorem{theorem}{Theorem}
\newtheorem{lemma}{Lemma}
\newtheorem{corollary}{Corollary}
\hyphenation{op-tical net-works semi-conduc-tor IEEE-Xplore}
% updated with editorial comments 8/9/2021

\begin{document}

\title{TopoStake: Accountable Propagation Incentives with Real-Stake-Bounded Security for PoS Blockchains}

\author{IEEE Publication Technology,~\IEEEmembership{Staff,~IEEE,}
        % <-this % stops a space
\thanks{This paper was produced by the IEEE Publication Technology Group. They are in Piscataway, NJ.}% <-this % stops a space
\thanks{Manuscript received April 19, 2021; revised August 16, 2021.}}

% The paper headers
\markboth{Journal of \LaTeX\ Class Files,~Vol.~14, No.~8, August~2021}%
{Shell \MakeLowercase{\textit{et al.}}: A Sample Article Using IEEEtran.cls for IEEE Journals}

% \IEEEpubid{0000--0000/00\$00.00~\copyright~2021 IEEE}
% Remember, if you use this you must call \IEEEpubidadjcol in the second
% column for its text to clear the IEEEpubid mark.

\maketitle

\begin{abstract}
This document describes the most common article elements and how to use the IEEEtran class with \LaTeX \ to produce files that are suitable for submission to the IEEE.  IEEEtran can produce conference, journal, and technical note (correspondence) papers with a suitable choice of class options. 
\end{abstract}

\begin{IEEEkeywords}
Article submission, IEEE, IEEEtran, journal, \LaTeX, paper, template, typesetting.
\end{IEEEkeywords}

\section{Introduction}
\label{sec:intro}

\IEEEPARstart{P}{roof-of-Stake} (PoS) blockchains secure ledger agreement by anchoring block production and finality voting in validators' economic stake~\cite{kiayias2017ouroboros,gilad2017algorand,buterin2017casper}.
However, this stake-based security still relies on a network service that is usually outside the consensus accounting: transactions, blocks, and votes must be propagated among validators in a timely manner before the protocol can make progress.
Propagation consumes bandwidth, computation, and verification effort, but relay behavior is often neither accountable nor directly incentivized at the protocol level.
As a result, fast dissemination benefits the whole validator set, yet the validators that provide relay service receive limited protocol-level recognition.
A propagation-aware PoS design must therefore reward relay contribution in an accountable way, while ensuring that relay-derived scores cannot weaken the real-stake basis of consensus security.

The absence of accountable relay incentives therefore creates a mismatch between stake-based rewards and relay behavior.
When relay effort remains unrewarded, rational validators may free-ride by reducing forwarding activity and relying on others to disseminate transactions and blocks~\cite{babaioff2012bitcoin,ersoy2017transaction}.
This behavior is individually reasonable because propagation consumes bandwidth and verification resources, while protocol rewards are usually tied to block proposal or voting rather than relay service~\cite{naumenko2019erlay,kang2019incentivizing}.
As free-riding spreads, the propagation burden shifts to fewer active nodes, reducing redundant dissemination paths and making the system more sensitive to delay, churn, and partial connectivity failures.
Prior studies have shown that propagation delay affects forks, confirmation latency, and the security-performance tradeoff of blockchain protocols~\cite{decker2013information,croman2016scaling,gervais2016security,pass2017analysis}.

Securely rewarding propagation requires more than attaching relay evidence to transactions and distributing fees~\cite{dai2026txtail}.
In a PoS system, such contribution becomes security-sensitive once it affects future proposer opportunities or other consensus-level incentives.
Relay paths may be selectively reported by proposers, padded by colluding relayers, or amplified through Sybil identities and transaction flooding.
A low-stake validator with a favorable network position may also obtain disproportionate propagation credit.
If such credit is allowed to directly increase finality voting power or unbounded proposer influence, the protocol would weaken the real-stake basis on which PoS security relies.

Existing work addresses related aspects of this problem, but still leaves open the combination of accountability, incentive compatibility, and PoS security preservation.
PoS and hybrid consensus protocols improve leader election, finality, or multi-resource robustness, but they usually treat network propagation as an underlying service rather than an accountable input to protocol incentives~\cite{liu2024bulwark,fitzi2022minotaur}.
Dissemination and mempool-accountability protocols reduce propagation delay, improve fairness, or make transaction handling auditable, but they do not connect relay behavior to consensus-level incentives~\cite{nasrulin2023lo,yahyaoui2025hermes,dai2025lowlatency}.
Transaction relay incentive schemes reward relayers using fee allocation and verifiable relay evidence, yet they mainly operate at the level of per-transaction dissemination and do not study how relay-derived contribution should affect PoS proposer selection under a real-stake security boundary~\cite{wang2023dualauction}.


In this paper, we propose \emph{TopoStake}, a propagation-accountability and incentive layer for PoS blockchains.
TopoStake is built around a separation between three quantities: real stake, propagation score, and bounded proposer-selection weight.
Real stake remains the basis of finality voting and is not increased by relay behavior.
Propagation score is computed from verifiable relay evidence and is used to allocate relay rewards.
To provide a long-term incentive for active dissemination, TopoStake further allows propagation score to affect proposer selection only through a capped bonus, so that relay-derived influence remains bounded by protocol parameters.
Through this design, TopoStake makes propagation contribution accountable and rewardable, while keeping consensus safety anchored in real stake.


\section{Preliminaries, Threat Model, and Design Goals}
\label{sec:system_model}

\subsection{Preliminaries}
\label{subsec:preliminaries}

\subsubsection{Network and Epochs}

We consider a permissionless PoS blockchain network.
During epoch $e$, the active validator set is
$\mathcal{V}_e=\{v_1,\dots,v_{N_e}\}$.
The communication network at slot $t$ is represented by a directed graph
$G_t=(\mathcal{V}_e,\mathcal{E}_t)$, where
$(v_i,v_j)\in\mathcal{E}_t$ means that $v_i$ can send messages to $v_j$.
The reverse direction is not implied.

Time is divided into fixed-duration slots and epochs.
The $e$-th epoch contains $L$ consecutive slots:
\begin{equation}
\mathcal{I}_e
=
\{t\mid (e-1)L\leq t<eL\}.
\end{equation}
For clarity, the validator set and the economic-stake snapshot are fixed within each epoch.
Stake changes take effect only at epoch boundaries according to the rules of the underlying PoS protocol.

The network is partially synchronous.
Before the unknown Global Stabilization Time (GST), message delay may be unbounded.
After GST, messages exchanged over honest communication paths are delivered within a known bound $\Delta$.
We further assume that the subgraph induced by honest validators remains connected after GST, so that honest validators are not permanently partitioned.

Validators use traceable gossip for transaction propagation.
Each relay verifies the existing path evidence and signs its incoming and outgoing propagation edges.
Section~\ref{subsec:path-tracing} gives the complete path construction.

\subsubsection{Epoch Snapshots and Protocol Weights}

Let $S_i(e)>0$ be the economic stake of validator $v_i$ in the snapshot used for epoch $e$.
Its normalized economic stake is
\begin{equation}
\label{eq:normalized_economic_stake}
\hat S_i(e)
=
\frac{S_i(e)}
{\sum_{v_j\in\mathcal{V}_e}S_j(e)}.
\end{equation}
Economic stake is the root security resource of TopoStake.
It is used directly for finality voting.

Propagation behavior observed in epoch $e$ is summarized by
$\mathrm{Score}_i(e)$ after the epoch is finalized.
This score is bound to the registered validator identity and is not transferable.
A newly registered validator starts with zero propagation score.

The score from epoch $e-1$ is used together with the economic-stake snapshot of epoch $e$ to derive the proposer-selection weight $W_i(e)$.
In contrast, relay fees are calculated from the verified path contribution in the containing block and are settled after that block is finalized.

TopoStake therefore separates three quantities:
economic stake $\hat S_i(e)$, propagation score $\mathrm{Score}_i(e)$, and proposer-selection weight $W_i(e)$.
Only economic stake is used for finality voting.

\subsubsection{Underlying PoS Protocol}

TopoStake is layered on an existing PoS protocol with a chain-selection rule, a randomness beacon, and a stake-weighted finality gadget.
We make the following modular assumptions about the underlying protocol.

First, if the adversarial proposer-selection share is below one half and the randomness beacon is unpredictable and sufficiently unbiased, the protocol provides probabilistic persistence, chain growth, and chain quality after GST.

Second, a block is finalized only after receiving votes representing at least two thirds of the economic-stake snapshot.
If Byzantine validators control less than one third of that stake, conflicting blocks cannot both be finalized.

Third, propagation metadata affects only relay rewards, propagation scores, and proposer bonuses.
A missing or invalid path record gives zero relay credit.
It does not invalidate the underlying transaction or stop the base PoS protocol from making progress.
Validators recompute relay payments from verified path records.

These assumptions allow the security analysis to separate propagation accountability from the safety of the underlying ledger.

\subsubsection{Security Properties}
\label{subsec:security_def}

Let $\mathcal{C}_t^i$ be the local chain of honest validator $v_i$ at slot $t$,
and let $\mathcal{C}_t^i\lceil k$ denote the chain after removing its last $k$ blocks.
Let $\mathcal{F}_t^i$ be the set of blocks finalized by $v_i$.

\begin{enumerate}
\item \textit{Safety.}
For any honest validators $v_i$ and $v_j$, the probability that a $k$-deep prefix of one honest chain is absent from another honest chain decreases exponentially in $k$:
\begin{equation}
\Pr\!\left[
\mathcal{C}_t^i\lceil k
\not\preceq
\mathcal{C}_{t'}^j
\right]
\leq
e^{-\Omega(k)}.
\end{equation}
In addition, honest validators do not finalize conflicting blocks at the same height.
If conflicting finality occurs, the signatures identify the validators that voted for both blocks.

\item \textit{Liveness.}
After GST, the ledger satisfies chain growth and chain quality.
There exist constants $\tau,\mu\in(0,1]$ and $s_0\in\mathbb{N}$ such that, for any honest validator and any window of $s\geq s_0$ slots,
\begin{equation}
|\mathcal{C}_{t+s}^i|-|\mathcal{C}_t^i|
\geq
\tau s,
\end{equation}
and at least a $\mu$ fraction of the blocks in the window are honestly proposed.
The value of $\mu$ depends on the honest share of proposer-selection weight.
\end{enumerate}

In addition to ledger safety and liveness, we analyze path authenticity, reward-budget balance, Sybil non-amplification, and bounded proposer influence.

\subsubsection{Cryptographic Assumptions}

TopoStake uses BLS signatures to compress the signatures on a propagation path.
Given signatures $\{\sigma_1,\dots,\sigma_n\}$ on messages
$\{m_1,\dots,m_n\}$, the aggregate signature is
\begin{equation}
\sigma_{\mathrm{agg}}
=
\sum_{i=1}^{n}\sigma_i.
\end{equation}
Its validity is checked by
\begin{equation}
e(\sigma_{\mathrm{agg}},g_2)
=
\prod_{i=1}^{n}
e\!\left(
H_{\mathbb{G}_1}(m_i),pk_i
\right).
\end{equation}

We assume existential unforgeability of the signature scheme,
collision resistance of the hash function,
and proofs of possession for registered BLS public keys.
The aggregate signature has constant size, but the validator sequence and verification work remain linear in the path length.


\subsection{Threat Model}
\label{subsec:problem_threat}

A computationally bounded adversary $\mathcal{A}$ may adaptively corrupt validators and control their forwarding behavior, path signatures, block proposals, and finality votes.
Let $\mathcal{A}_e\subseteq\mathcal{V}_e$ be the corrupted validators during epoch $e$.
Their economic-stake fraction is
\begin{equation}
\label{eq:adversarial_stake_fraction}
f_{\mathcal{A}}^{S}(e)
=
\sum_{v_i\in\mathcal{A}_e}
\hat S_i(e).
\end{equation}
For finality safety, we assume
\begin{equation}
\label{eq:real_stake_assumption}
f_{\mathcal{A}}^{S}(e)
<
\frac{1}{3}
\end{equation}
throughout each finality voting period.

The adversary's network position is not assumed to be proportional to its stake.
Corrupted validators may have high degree, large bandwidth, or advantageous positions in the propagation graph.
They may also provide apparently valid relay service and obtain high propagation scores.
Such scores do not increase their finality voting weight.

Let
\begin{equation}
\label{eq:adversarial_proposer_fraction}
f_{\mathcal{A}}^{W}(e)
=
\frac{
\sum_{v_i\in\mathcal{A}_e}W_i(e)
}{
\sum_{v_j\in\mathcal{V}_e}W_j(e)
}
\end{equation}
be the adversarial proposer-selection share.
Section~\ref{sec:security} proves that protocol parameters can ensure
$f_{\mathcal{A}}^{W}(e)<1/2$ whenever
$f_{\mathcal{A}}^{S}(e)<1/3$.

We consider the following attacks:

\begin{enumerate}
\item \textit{Path forgery:}
the adversary attempts to insert, remove, or reorder honest validators in a propagation path.

\item \textit{Receipt equivocation:}
a validator signs several incompatible incoming paths for the same transaction and epoch.

\item \textit{Sybil splitting and path padding:}
the adversary divides stake across identities or inserts consecutive controlled identities into a path.

\item \textit{Transaction flooding:}
the adversary creates transactions to increase relay rewards or propagation scores.

\item \textit{Proposer manipulation:}
a proposer selects valid paths that favor colluding relayers or omits paths that favor honest relayers.

\item \textit{Lazy propagation:}
rational validators reduce forwarding effort while attempting to retain normal PoS rewards.
\end{enumerate}

A fully colluding set of validators can sign a syntactically valid path.
TopoStake does not claim that signatures prove physical latency or useful network service.
Instead, it prevents the fabrication of honest participation and bounds the reward and proposer influence obtainable from any accepted path.

The adversary may delay, drop, or reorder messages passing through corrupted validators.
After GST, it cannot prevent delivery over honest communication paths.
Security proofs treat non-corrupted validators as protocol-following.
The economic effect of rational relay behavior is evaluated separately from ledger safety.


\subsection{Design Goals}
\label{subsec:goals}

\begin{figure}[htbp]
  \centering
  \includegraphics[width=\linewidth]{figures/topology_comparison.pdf}
  \caption{Propagation-unaware and propagation-accountable incentives.}
  \label{fig:motivation}
\end{figure}

PoS protocols reward block production and finality voting, but usually do not reward transaction relay.
This creates a free-rider problem: validators can rely on other nodes to propagate transactions while avoiding bandwidth and computation costs.
It can also leave active low-stake relayers with little economic recognition.

TopoStake adds a propagation-accountability layer without replacing economic stake as the security resource.
Its design has four goals:

\begin{enumerate}
\item \textbf{Verifiable path evidence.}
Relay credit must be based on publicly verifiable path records rather than self-reported claims.

\item \textbf{Budgeted and accountable rewards.}
Each transaction must have a bounded relay budget.
Invalid paths, repeated identities, and conflicting receipts must not create relay credit.

\item \textbf{Real-stake-bounded proposer influence.}
Propagation scores may improve proposer opportunities, but the increase must be capped by protocol parameters.
Finality voting must remain based only on economic stake.

\item \textbf{Improved relay participation.}
Immediate relay fees and bounded future proposer opportunities should encourage sustained forwarding behavior.
Topology improvement is treated as an empirical result rather than an unconditional security guarantee.
\end{enumerate}


\section{The TopoStake Protocol}
\label{sec:protocol}

\subsection{Overview}

TopoStake adds a propagation-accountability layer to a PoS blockchain.
It makes relay behavior verifiable and rewardable, while keeping finality voting based on real economic stake.
As shown in Fig.~\ref{fig:overview}, TopoStake contains four stages:
\emph{Path Tracing}, \emph{Contribution Quantification},
\emph{Bounded Proposer Selection}, and \emph{Reward Allocation}.

\textbf{Step 1: Path Tracing.}
Each relay signs the incoming and outgoing propagation edges of a transaction.
The resulting path is publicly verifiable and is compressed into one BLS aggregate signature for block inclusion.

\textbf{Step 2: Contribution Quantification.}
Each verified path receives a bounded contribution budget.
The budget is divided among intermediate relayers and decreases with path length.
Path contributions are then aggregated and smoothed into a long-term propagation score.

\textbf{Step 3: Bounded Proposer Selection.}
Economic stake remains the base proposer-selection weight.
A propagation score provides only a capped bonus, so proposer influence remains bounded by real stake and protocol parameters.
Finality voting does not use this bonus.

\textbf{Step 4: Reward Allocation.}
A fraction of each transaction fee is paid to intermediate relayers according to their path-level contribution.
The same contributions update future propagation scores, creating both immediate relay rewards and bounded future proposer opportunities.


\subsection{Chained-Signature Path Tracing}
\label{subsec:path-tracing}

\subsubsection{Two-sided path certificates}
Consider a propagation path
$P=(v_0,v_1,\dots,v_m)$,
where $v_0$ is the transaction originator and $v_m$ is the block proposer.
Each transaction carries a route array
$R=[r_0,\dots,r_{m-1}]$.

Let $e$ be the current epoch and let $H(\cdot)$ be a collision-resistant hash function.
The originator initializes
\begin{equation}
\label{eq:path_chain_init}
c_0 =
H\bigl(H(tx)\mathbin\Vert e\mathbin\Vert H(v_0)\bigr).
\end{equation}
For $i>0$, the chain value is updated as
\begin{equation}
\label{eq:path_chain_update}
c_i =
H\bigl(c_{i-1}\mathbin\Vert H(v_i)\bigr).
\end{equation}
Thus, $c_i$ commits to the transaction, the epoch, and the ordered path prefix ending at $v_i$.

For the edge $(v_i,v_{i+1})$, define
\begin{equation}
\label{eq:edge_statement}
M_i =
c_i\mathbin\Vert H(v_i)\mathbin\Vert H(v_{i+1}).
\end{equation}
The sender and receiver proofs are
\begin{equation}
\label{eq:path_proofs}
\begin{aligned}
\sigma_i^{\mathsf{S}}
&=
\mathsf{Sign}(sk_i,M_i),
&& 0\leq i<m,\\
\sigma_i^{\mathsf{R}}
&=
\mathsf{Sign}(sk_i,M_{i-1}),
&& 0<i\leq m.
\end{aligned}
\end{equation}

The route array uses a node-centric encoding:
\begin{equation}
\label{eq:route_entries}
\begin{aligned}
r_0
&=
\bigl(v_1,\sigma_0^{\mathsf{S}}\bigr),\\
r_i
&=
\bigl(v_{i+1},\sigma_i^{\mathsf{R}},
      \sigma_i^{\mathsf{S}}\bigr),
&& 0<i<m.
\end{aligned}
\end{equation}
The proposer does not create a route entry because it has no next hop.
Its receiving proof $\sigma_m^{\mathsf{R}}$ is included in the final aggregate proof.

A hop $(v_i,v_{i+1})$ is valid only if
\begin{align}
\mathsf{Verify}\!\left(
pk_i,M_i,\sigma_i^{\mathsf{S}}
\right)
&=\mathsf{true}, \\
\mathsf{Verify}\!\left(
pk_{i+1},M_i,\sigma_{i+1}^{\mathsf{R}}
\right)
&=\mathsf{true}.
\end{align}
The proposer verifies all hop certificates before including the path in a block.
Paths containing repeated validator identities are rejected.

For each transaction and epoch, a validator signs at most one receiving proof.
This rule commits the validator to one incoming predecessor for contribution accounting.
It does not restrict gossip fanout: the validator may still forward the transaction to several peers and create sender proofs for several outgoing edges.
Two receiving proofs signed by the same validator for the same transaction and epoch form public equivocation evidence.

\subsubsection{Path selection}
At most one path is recorded for each transaction in a block.
An honest proposer selects the shortest valid path in its local cache and uses a canonical ordering to break ties.
This is a local rule rather than a global fastest-path guarantee.
A verifier can check path validity, but cannot prove that the proposer did not observe another path.
The security analysis therefore bounds the effect of path selection instead of assuming an honest proposer.

\subsubsection{BLS compression and block encoding}
For a path $P=(v_0,\dots,v_m)$, the proposer aggregates all sender and receiver proofs as
\begin{equation}
\label{eq:path_aggregate}
\sigma_{\mathrm{agg}}
=
\sigma_0^{\mathsf{S}}
+
\sum_{i=1}^{m-1}
\left(
\sigma_i^{\mathsf{R}}+\sigma_i^{\mathsf{S}}
\right)
+
\sigma_m^{\mathsf{R}}
\in\mathbb{G}_1.
\end{equation}
The aggregate is verified using the standard BLS aggregate-verification equation over the corresponding $2m$ signer-message pairs.
The aggregate signature has constant size, although verification and the validator sequence remain linear in $m$.

A block is represented as
\begin{equation}
\label{eq:block_struct}
\mathcal{B}
=
(\mathcal{H},\mathcal{T},\mathcal{P}),
\end{equation}
where $\mathcal{H}$ is the block header,
$\mathcal{T}=\{tx_0,\dots,tx_n\}$ is the transaction set, and
$\mathcal{P}=\{p_0,\dots,p_n\}$ contains the path records.

\begin{figure}[htbp]
  \centering
  \includegraphics[width=\linewidth]{figures/block_structure.pdf}
  \caption{Block-level path encoding in TopoStake.}
  \label{fig:block}
\end{figure}

For transaction $tx_k$, the path record is
\begin{equation}
\label{eq:path_record}
p_k =
\left(
\sigma_{\mathrm{agg}}^k,
[v_0^k,v_1^k,\dots,v_{m_k-1}^k]
\right).
\end{equation}
The proposer $v_{m_k}^k$ is obtained from the block header and is therefore not repeated in the sequence.
Given the transaction, the stored sequence, and the proposer identity, a verifier can reconstruct all edge statements and verify the aggregate proof.

If the proposer generated the transaction, the path record is empty.
If the originator sends the transaction directly to the proposer, the path contains no intermediate relayer and receives no relay contribution.

The certificate proves that both endpoints endorsed each recorded edge.
It does not prove physical latency or prevent a fully colluding path.
Such paths are instead limited by the path budget, score saturation, and bounded proposer influence defined below.


\subsection{Quantification of Propagation Contribution}
\label{subsec:contribution-quantification}

TopoStake scores only paths that pass public verification.
Each transaction contributes at most one path, and only intermediate nodes receive relay credit.
For a path $p=(v_0,\dots,v_m)$, define
\begin{equation}
\label{eq:relay_set}
\mathcal{R}(p)
=
\{v_1,\dots,v_{m-1}\}.
\end{equation}
The originator and proposer are excluded.
A direct path with $m\leq1$ has no eligible relayer.

\subsubsection{Target depth and path budget}
TopoStake maintains a target propagation depth $D_e\geq1$ for each epoch.
Let $\widetilde L_e$ be the median hop length of eligible paths finalized in epoch $e$.
The target changes by at most one hop per epoch:
\begin{equation}
\label{eq:target_update}
D_{e+1}
=
\begin{cases}
D_e, & \mathcal{P}_e=\emptyset,\\
D_e+1, & \widetilde L_e>D_e,\\
D_e, & \widetilde L_e=D_e,\\
\max\{1,D_e-1\}, & \widetilde L_e<D_e.
\end{cases}
\end{equation}
Using the median and a one-step update reduces the effect of short-term path manipulation.
The value of $D_e$ changes reward scaling only; it does not affect finality voting.

For a path of $m\geq2$ hops, define its total contribution budget as
\begin{equation}
\label{eq:path_budget_def}
B_e(m)
=
\min\left\{1,\frac{D_e}{m}\right\}
\lambda_e^{m-1},
\qquad
\lambda_e
=
\frac{2D_e+1}{3D_e+1}.
\end{equation}
Since $0<\lambda_e<1$, longer paths receive a smaller budget.

The budget is divided among the intermediate relayers using
\begin{equation}
\label{eq:position_weight}
\alpha_k^{(m,e)}
=
\frac{(1-r_e)r_e^{k-1}}
{1-r_e^{m-1}},
\qquad
r_e
=
\frac{D_e}{2D_e+1},
\end{equation}
where $1\leq k\leq m-1$.
The contribution of relay $v_k$ is
\begin{equation}
\label{eq:atomic_credit}
\gamma(v_k,p)
=
B_e(m)\alpha_k^{(m,e)}.
\end{equation}

The construction satisfies


\begin{align}
\sum_{k=1}^{m-1}\alpha_k^{(m,e)} &= 1, \label{eq:path_credit_1} \\
\sum_{v_k\in\mathcal{R}(p)} \gamma(v_k,p) &= B_e(m)\leq1, \label{eq:path_credit_2} \\
\lambda_e(1+r_e) &= 1. \label{eq:path_credit_3}
\end{align}
The first two identities give every path a bounded budget.
The last identity is used in Section~\ref{sec:security} to show that replacing one relay position with consecutive Sybil identities does not increase their total path credit.
The decreasing positional weights reward nodes that make the transaction available earlier along the recorded path.

\subsubsection{Epoch-level score update}
Let $\mathcal{P}_e(v_i)$ be the eligible paths finalized in epoch $e$ that contain $v_i$ as an intermediate relayer.
The raw contribution of $v_i$ is
\begin{equation}
\label{eq:raw_epoch_credit}
A_i(e)
=
\sum_{p\in\mathcal{P}_e(v_i)}
\gamma(v_i,p).
\end{equation}

TopoStake applies a stake-scaled logarithmic function:
\begin{equation}
\label{eq:saturated_credit}
C_i(e)
=
\hat S_i(e)
\log\left(
1+
\frac{A_i(e)}
{K\hat S_i(e)}
\right),
\end{equation}
where $K>0$ is a saturation parameter.
The logarithm gives diminishing returns to repeated relay activity.
Its stake-scaled form also supports the stake-splitting bound proved in Section~\ref{sec:security}.

The long-term propagation score is updated by EMA:
\begin{equation}
\label{eq:score_ema}
\mathrm{Score}_i(e)
=
\beta C_i(e)
+
(1-\beta)\mathrm{Score}_i(e-1),
\end{equation}
where $\beta\in(0,1)$ and $\mathrm{Score}_i(0)=0$.
The score is computed after epoch $e$ is finalized and is used only from epoch $e+1$.

For proposer selection, scores are normalized as
\begin{equation}
\label{eq:normalized_score}
\hat C_i(e)
=
\begin{cases}
\dfrac{\mathrm{Score}_i(e)}
{\sum_{v_j\in\mathcal{V}_e}\mathrm{Score}_j(e)},
&
\sum_{v_j\in\mathcal{V}_e}\mathrm{Score}_j(e)>0,\\[1.2ex]
0, & \text{otherwise}.
\end{cases}
\end{equation}


\subsection{Bounded Influence Weight and Leader Election}
\label{subsec:bounded-influence}

TopoStake uses propagation scores only for proposer selection.
The underlying finality gadget continues to use economic stake.

\subsubsection{Bounded proposer weight}
A validator receives a propagation bonus in epoch $e$ only when its previous propagation-score share exceeds its current stake share:
\begin{equation}
\label{eq:prop_bonus}
b_i(e)
=
\min\left\{
b_{\max},
\left[
\frac{\hat C_i(e-1)}
{\hat S_i(e)}
-1
\right]^+
\right\},
\end{equation}
where $[x]^+=\max\{x,0\}$ and $b_{\max}$ is the maximum bonus.

The proposer-selection weight for epoch $e$ is
\begin{equation}
\label{eq:bounded_weight}
W_i(e)
=
\hat S_i(e)
\bigl(1+\eta b_i(e)\bigr),
\end{equation}
where $\eta\geq0$ controls the strength of the propagation incentive.
The normalized weight is
\begin{equation}
\label{eq:normalized_weight}
\hat W_i(e)
=
\frac{W_i(e)}
{\sum_{v_j\in\mathcal{V}_e}W_j(e)}.
\end{equation}

Because $b_i(e)\leq b_{\max}$,
\begin{equation}
\label{eq:weight_bound}
W_i(e)
\leq
\bigl(1+\eta b_{\max}\bigr)\hat S_i(e).
\end{equation}
For any adversarial coalition $\mathcal{A}$, this implies
\begin{equation}
\label{eq:coalition_weight_bound}
f_{\mathcal{A}}^{W}(e)
\leq
\bigl(1+\eta b_{\max}\bigr)
f_{\mathcal{A}}^{S}(e).
\end{equation}
TopoStake requires
\begin{equation}
\label{eq:bonus_parameter_bound}
\eta b_{\max}\leq\frac{1}{2}.
\end{equation}
Therefore, an adversary with less than one third of the real stake has less than one half of the proposer-selection weight.
When no previous propagation score exists, all bonuses are zero and proposer selection reduces to ordinary PoS.


\subsubsection{Leader election and finality}
TopoStake reuses the unpredictable randomness beacon of the underlying PoS protocol, such as a bias-resistant RANDAO or VRF construction.
Within epoch $e+1$, validator $v_i$ is selected as proposer with probability
\begin{equation}
\label{eq:leader_probability}
\Pr[v_i\text{ is selected}]
=
\hat W_i(e).
\end{equation}

Finality voting uses only the economic-stake snapshot:
\begin{equation}
\label{eq:finality_weight}
V_i^{\mathrm{final}}(e)
=
\hat S_i(e).
\end{equation}
A finality quorum $Q$ must satisfy
\begin{equation}
\label{eq:finality_quorum}
\sum_{v_i\in Q}\hat S_i(e)
\geq
\frac{2}{3}.
\end{equation}
Propagation scores and proposer bonuses therefore affect block-production opportunities, but not finality voting power.


\subsection{Incentive and Reward Mechanism}
\label{subsec:reward}

TopoStake provides two incentives for relay service:
an immediate share of transaction fees and a bounded increase in future proposer opportunities.

\subsubsection{Fee split and relay reward}
Let $f(tx)$ be the fee of transaction $tx$, and let $\theta\in[0,1)$ be the proposer fee ratio.
For block $\mathcal{B}$, the proposer receives
\begin{equation}
\label{eq:proposer_reward}
R_{\mathrm{prop}}(\mathcal{B})
=
R_{\mathrm{base}}(\mathcal{B})
+
\theta
\sum_{tx\in\mathcal{T}}f(tx).
\end{equation}

For an eligible path $p$ of transaction $tx$, intermediate relay $v_i$ receives
\begin{equation}
\label{eq:relay_reward}
R_i^{\mathrm{relay}}(tx)
=
(1-\theta)f(tx)\gamma(v_i,p).
\end{equation}
Because the total path credit is at most one,
\begin{align}
\label{eq:relay_reward_bound}
\sum_{v_i\in\mathcal{R}(p)}
R_i^{\mathrm{relay}}(tx)
&=
(1-\theta)f(tx)B_e(m) \notag\\
&\leq
(1-\theta)f(tx).
\end{align}
Thus, relay payments never exceed the relay-fee budget.
Long or padded paths receive a smaller total payment because $B_e(m)$ decreases with $m$.
Any undistributed part of the relay budget is burned and cannot be claimed by the proposer or the relayers on that path.

Relay rewards are settled only after the containing block is finalized.
The originator and proposer receive no relay reward.
A proposer-generated transaction or a direct originator-to-proposer path therefore produces no relay payment.

The same values $\gamma(v_i,p)$ update the epoch-level propagation score.
Current paths determine current relay payments, while finalized scores affect only future proposer weights.
This avoids a circular dependence within the same block or epoch.

\subsubsection{Invalid and conflicting evidence}
An invalid path receives no relay reward and contributes no score.
A path is invalid if a hop certificate fails verification, the path repeats a validator identity, or the path record is malformed.

If validator $v_i$ signs two different receiving proofs for the same transaction and epoch, both signatures form public equivocation evidence.
The conflicting paths are excluded, and the validator's propagation score for the next epoch is revoked:
\begin{equation}
\label{eq:score_revocation}
\mathrm{Score}_i(e+1)=0.
\end{equation}
This penalty affects relay rewards and proposer bonuses only.
Finality equivocation remains subject to the real-stake slashing rules of the underlying PoS protocol.



\section{Security Analysis}
\label{sec:security}

This section analyzes the security properties introduced in Sections~\ref{sec:system_model} and~\ref{sec:protocol}.
The analysis is modular.
TopoStake proves the authenticity and bounded economic effect of propagation evidence, while ledger safety and liveness are inherited from the underlying PoS protocol under the proposer- and stake-threshold conditions stated earlier.

\subsection{Path Authenticity and Accountability}

\begin{theorem}[Path Authenticity]
\label{thm:path_authenticity}
Suppose a propagation path is accepted by an honest verifier.
Except with negligible probability, every honest validator recorded as an endpoint of a hop signed the corresponding edge statement, and the accepted validator sequence is bound to the transaction and epoch.
\end{theorem}

\begin{proof}
For hop $(v_i,v_{i+1})$, a valid certificate contains a sender signature from $v_i$ and a receiver signature from $v_{i+1}$ on the same edge statement $M_i$.
If either endpoint is honest, producing its signature without its participation would break signature unforgeability.

The initial chain value contains the transaction hash, epoch, and originator identity.
Each later chain value contains the previous chain value and the next validator identity.
Therefore, inserting, removing, substituting, or reordering an earlier validator changes the downstream edge statements.
Keeping the old signatures valid would require either a signature forgery or a hash collision.
\end{proof}

Theorem~\ref{thm:path_authenticity} proves endorsement, not physical transmission time.
If both endpoints are corrupted, they may jointly endorse an edge.
The later reward and influence bounds limit the effect of such collusion.

\begin{lemma}[Receipt Accountability]
\label{lem:receipt_accountability}
If a validator signs two receiving proofs for the same transaction and epoch but with different predecessors or path prefixes, the two signatures form publicly verifiable equivocation evidence.
An honest validator cannot be falsely accused except with negligible probability.
\end{lemma}

\begin{proof}
The two receiving proofs are signatures under the same validator key on two different edge statements.
Both statements are bound to the same transaction and epoch through their chain values.
Any verifier can check both signatures and compare the predecessor or prefix.
Creating such evidence without the validator's signatures would require a forgery.
\end{proof}

The single-receipt rule does not restrict gossip fanout.
A validator may forward a transaction to several peers, but only one incoming edge is eligible for contribution accounting.


\subsection{Reward-Budget and Padding Bounds}

\begin{theorem}[Path-Budget Balance]
\label{thm:path_budget}
For every eligible path $p$ with $m\geq2$ hops,
\begin{equation}
\sum_{v_i\in\mathcal{R}(p)}
\gamma(v_i,p)
=
B_e(m)
\leq
1.
\end{equation}
Moreover, for fixed $D_e$, $B_e(m)$ decreases as $m$ increases.
Consequently, the total relay payment for transaction $tx$ satisfies
\begin{equation}
\sum_{v_i\in\mathcal{R}(p)}
R_i^{\mathrm{relay}}(tx)
\leq
(1-\theta)f(tx).
\end{equation}
\end{theorem}

\begin{proof}
The positional weights are normalized:
\begin{equation}
\sum_{k=1}^{m-1}
\alpha_k^{(m,e)}
=
1.
\end{equation}
Therefore, summing
$\gamma(v_k,p)=B_e(m)\alpha_k^{(m,e)}$
over all intermediate relayers gives $B_e(m)$.

The efficiency term
$\min\{1,D_e/m\}$ is non-increasing in $m$,
and $0<\lambda_e<1$.
Hence
\begin{equation}
B_e(m)
=
\min\left\{1,\frac{D_e}{m}\right\}
\lambda_e^{m-1}
\end{equation}
decreases with path length.
Multiplying the path budget by $(1-\theta)f(tx)$ gives the relay-payment bound.
\end{proof}

This theorem means that a proposer may include a long valid path, but a longer path cannot create a larger total relay pool.
Any undistributed fee is burned and cannot be returned to the proposer or the selected relayers.

\begin{theorem}[Consecutive Identity-Padding Non-Amplification]
\label{thm:path_padding}
Consider a relay at position $k$ in an $m$-hop path.
Suppose the relay is replaced by $q+1$ consecutive identities controlled by the same coalition, increasing the path length to $m+q$.
For fixed $D_e$, the total atomic credit of the inserted identities is no greater than the credit of the original relay.
\end{theorem}

\begin{proof}
Write
\begin{equation}
\psi_e(m)
=
\min\left\{1,\frac{D_e}{m}\right\}.
\end{equation}
The original relay credit is
\begin{equation}
\Gamma_0
=
\psi_e(m)\lambda_e^{m-1}
\frac{(1-r_e)r_e^{k-1}}
{1-r_e^{m-1}}.
\end{equation}
After inserting $q$ additional identities, their total credit is
\begin{equation}
\Gamma_q
=
\psi_e(m+q)\lambda_e^{m+q-1}
\frac{(1-r_e)r_e^{k-1}
\sum_{\ell=0}^{q}r_e^\ell}
{1-r_e^{m+q-1}}.
\end{equation}
Their ratio satisfies
\begin{equation}
\frac{\Gamma_q}{\Gamma_0}
=
\frac{\psi_e(m+q)}{\psi_e(m)}
\lambda_e^q
\left(
\sum_{\ell=0}^{q}r_e^\ell
\right)
\frac{1-r_e^{m-1}}
{1-r_e^{m+q-1}}.
\end{equation}
The first and last factors are at most one.
Since
$\lambda_e(1+r_e)=1$ and
\begin{equation}
\sum_{\ell=0}^{q}r_e^\ell
\leq
(1+r_e)^q,
\end{equation}
the middle factors are also at most one.
Thus $\Gamma_q\leq\Gamma_0$.
\end{proof}

Theorem~\ref{thm:path_padding} addresses identity splitting inside an existing relay position.
It does not prove that a proposer always selects an honest path.
A malicious proposer may favor a fully colluding valid path, but Theorem~\ref{thm:path_budget} still limits the total payment for that transaction.


\subsection{Score Robustness}

\begin{theorem}[Stake-Splitting Non-Amplification]
\label{thm:stake_splitting}
Suppose a coalition has total normalized stake $s$ and total raw contribution $a$ in an epoch.
If it splits these resources among identities
$(s_1,a_1),\dots,(s_n,a_n)$, where
\begin{equation}
\sum_{j=1}^{n}s_j=s,
\qquad
\sum_{j=1}^{n}a_j=a,
\end{equation}
then the aggregate saturated contribution of the split identities is no greater than that of one identity holding $(s,a)$:
\begin{equation}
\sum_{j=1}^{n}
s_j
\log\left(
1+\frac{a_j}{Ks_j}
\right)
\leq
s
\log\left(
1+\frac{a}{Ks}
\right).
\end{equation}
\end{theorem}

\begin{proof}
Let $w_j=s_j/s$.
Then $\sum_j w_j=1$.
By concavity of the logarithm,
\begin{align}
\sum_{j=1}^{n}
s_j
\log\left(
1+\frac{a_j}{Ks_j}
\right)
&=
s\sum_{j=1}^{n}
w_j
\log\left(
1+\frac{a_j}{Ks_j}
\right) \notag\\
&\leq
s\log\left(
1+
\sum_{j=1}^{n}
w_j\frac{a_j}{Ks_j}
\right) \notag\\
&=
s\log\left(
1+\frac{a}{Ks}
\right).
\end{align}
\end{proof}

The EMA update is a linear combination of the current saturated contribution and the previous score.
Therefore, it does not reverse the inequality in Theorem~\ref{thm:stake_splitting}.
New identities start with zero score, and existing scores cannot be transferred to newly created identities.

\begin{lemma}[Fee-Bounded Transaction Flooding]
\label{lem:flooding}
Consider a self-generated transaction with fee $f$.
Even if the adversary controls the transaction originator, proposer, and all intermediate relayers, the total fee returned to the adversarial coalition is at most $f$:
\begin{equation}
R_{\mathcal{A}}(tx)
\leq
\theta f
+
(1-\theta)fB_e(m)
\leq
f.
\end{equation}
If $\theta<1$ and the transaction contains an intermediate relay path, the returned amount is strictly less than $f$.
\end{lemma}

\begin{proof}
The proposer receives at most $\theta f$.
By Theorem~\ref{thm:path_budget}, all relayers together receive at most
$(1-\theta)fB_e(m)$.
Since $B_e(m)\leq1$, the total is at most $f$.
For $m\geq2$, $\lambda_e^{m-1}<1$, so $B_e(m)<1$ and part of the fee is burned.
\end{proof}

Thus, transaction flooding cannot mint a positive fee profit.
An adversary may still spend fees to purchase propagation score, but logarithmic saturation limits marginal score growth, and the proposer-weight cap below limits its consensus effect.


\subsection{Bounded Proposer Influence}

\begin{theorem}[Real-Stake-Bounded Proposer Influence]
\label{thm:bounded_influence}
For any adversarial coalition $\mathcal{A}$ in epoch $e$,
\begin{equation}
f_{\mathcal{A}}^{W}(e)
\leq
\bigl(1+\eta b_{\max}\bigr)
f_{\mathcal{A}}^{S}(e).
\end{equation}
\end{theorem}

\begin{proof}
For every validator,
$0\leq b_i(e)\leq b_{\max}$.
Therefore,
\begin{equation}
\sum_{v_i\in\mathcal{A}_e}
W_i(e)
\leq
\bigl(1+\eta b_{\max}\bigr)
\sum_{v_i\in\mathcal{A}_e}
\hat S_i(e).
\end{equation}
Since all bonuses are non-negative,
\begin{equation}
\sum_{v_j\in\mathcal{V}_e}W_j(e)
\geq
\sum_{v_j\in\mathcal{V}_e}\hat S_j(e)
=
1.
\end{equation}
Dividing the two expressions gives the result.
\end{proof}

\begin{corollary}[Honest Proposer Majority]
\label{cor:honest_proposer_majority}
If
\begin{equation}
f_{\mathcal{A}}^{S}(e)<\frac{1}{3}
\qquad\text{and}\qquad
\eta b_{\max}\leq\frac{1}{2},
\end{equation}
then
\begin{equation}
f_{\mathcal{A}}^{W}(e)<\frac{1}{2}.
\end{equation}
\end{corollary}

\begin{proof}
By Theorem~\ref{thm:bounded_influence},
\begin{equation}
f_{\mathcal{A}}^{W}(e)
\leq
\frac{3}{2}
f_{\mathcal{A}}^{S}(e)
<
\frac{1}{2}.
\end{equation}
\end{proof}

This bound holds regardless of how the adversary distributes its propagation score among identities.
Therefore, an advantageous network position may increase proposer opportunities, but cannot create an unbounded proposer majority.


\subsection{Ledger Safety and Liveness}

\begin{theorem}[Real-Stake Finality Safety]
\label{thm:finality_safety}
If Byzantine validators control less than one third of the economic-stake snapshot, two conflicting blocks cannot both be finalized by honest validators.
If conflicting finality occurs, the signed votes identify at least one third of the economic stake that voted for both blocks.
\end{theorem}

\begin{proof}
Let $Q_1$ and $Q_2$ be the voting quorums for two conflicting blocks.
Each quorum contains at least two thirds of the economic stake.
Therefore,
\begin{equation}
\hat S(Q_1\cap Q_2)
\geq
\hat S(Q_1)+\hat S(Q_2)-1
\geq
\frac{1}{3}.
\end{equation}
Every validator in the intersection signed both conflicting blocks.
Honest validators do not sign conflicting finality votes.
Hence conflicting finality would require at least one third of the economic stake to equivocate, contradicting
$f_{\mathcal{A}}^{S}<1/3$.
The two sets of signed votes also identify the equivocating validators.
\end{proof}

Propagation scores do not appear in this proof.
A validator may obtain a larger proposer probability, but its finality voting weight remains its economic-stake share.

\begin{theorem}[Persistence, Chain Growth, and Chain Quality]
\label{thm:ledger_liveness}
Assume GST has occurred, the honest communication graph satisfies the connectivity assumption, the randomness beacon satisfies the assumptions of the base PoS protocol, and the conditions of Corollary~\ref{cor:honest_proposer_majority} hold.
Then TopoStake inherits probabilistic persistence, chain growth, and chain quality from the underlying PoS protocol.
\end{theorem}

\begin{proof}
Corollary~\ref{cor:honest_proposer_majority} gives
$f_{\mathcal{A}}^{W}(e)<1/2$.
Thus, honest validators hold a strict majority of proposer-selection weight.

The randomness beacon selects proposers according to the normalized weights.
Over a sufficiently long slot window, standard concentration bounds imply a linear number of honest proposer slots except with negligible probability.
After GST, blocks and votes created by honest validators are delivered within the network bound.
The persistence, chain-growth, and chain-quality results of the underlying PoS protocol therefore apply.

Propagation records do not weaken this argument.
A missing or invalid record produces no reward or score, but does not invalidate the underlying block or prevent finality.
\end{proof}


\subsection{Security Scope and Limitations}

The above results establish the following boundaries.

First, TopoStake proves that honest validator identities cannot be inserted into a path without their signatures.
It does not prove that a fully corrupted path represents useful physical propagation.

Second, a malicious proposer may select a valid path that favors colluding relayers.
The protocol does not prove globally optimal path selection.
Instead, it bounds the total payment of every path and bounds all propagation-derived proposer influence by real stake.

Third, the target depth $D_e$ may affect the relative reward assigned to paths.
However, Theorems~\ref{thm:path_budget} and~\ref{thm:bounded_influence} hold for every $D_e\geq1$.
Therefore, manipulating $D_e$ may distort reward allocation, but it cannot violate the relay-fee budget, create unbounded proposer weight, or change finality voting power.

Finally, the security analysis does not claim that local timestamps prove global propagation latency.
Timing information is used only for epochs and scoring windows.


\section{Evaluation}
\label{sec:evaluation}

We evaluate TopoStake with five goals. First, we measure whether propagation accountability can be added to a PoS blockchain with moderate end-to-end overhead. Second, we quantify the cryptographic cost of two-sided BLS path certificates and aggregate verification. Third, we test whether the real-stake-bounded proposer influence derived in Section~\ref{sec:security} holds under adversaries with advantageous network positions. Fourth, we evaluate the robustness of the reward mechanism against path padding, transaction flooding, and heterogeneous relay behavior. Fifth, we isolate whether the persistent, bounded proposer bonus provides a dependability benefit beyond transaction-scoped relay fees under sustained validator outages.

\subsection{Experimental Setup}
\label{subsec:eval_setup}

We implement TopoStake in an event-driven Rust simulator. The simulator supports a base PoS protocol and the TopoStake extension, and generates peer-to-peer networks using Barabasi--Albert (BA), Erdős--Rényi (ER), and Watts--Strogatz (WS) graph models. Unless otherwise specified, we use BA as the default topology, repeat each configuration with three deterministic seeds, and report the median together with dispersion across seeds.

We compare three main configurations: \emph{PoS}, which uses economic stake for proposer selection and does not reward transaction relay; \emph{TopoStake-$\eta0$}, which enables relay rewards and propagation scoring but disables the proposer bonus; and \emph{TopoStake}, which enables both relay rewards and bounded proposer bonuses. This separation allows us to distinguish the immediate relay incentive from the bounded future proposer opportunity.

We record throughput, block success rate, transaction confirmation latency, path length, valid and invalid path records, relay rewards, proposer rewards, reward concentration, and adversarial proposer influence. For adversarial experiments, we additionally record the adversarial real-stake share, propagation-score share, proposer-weight share, observed proposer frequency, theoretical proposer-weight bound, and bound violations.

\subsection{Performance and Scalability}
\label{subsec:eval_performance}

\begin{figure*}[!t]
    \centering
    \subfloat[Throughput vs. input rate\label{fig:throughput_rate}]{
        \includegraphics[width=0.234\textwidth]{figures/performance_load_throughput.pdf}
    }
    \hfill
    \subfloat[p95 latency vs. input rate\label{fig:latency_rate}]{
        \includegraphics[width=0.234\textwidth]{figures/performance_load_latency.pdf}
    }
    \hfill
    \subfloat[Throughput vs. validators\label{fig:throughput_validators}]{
        \includegraphics[width=0.234\textwidth]{figures/performance_scale_throughput.pdf}
    }
    \hfill
    \subfloat[p95 latency vs. validators\label{fig:latency_validators}]{
        \includegraphics[width=0.234\textwidth]{figures/performance_scale_latency.pdf}
    }
    \caption{Performance of TopoStake under different transaction loads and network scales.}
    \label{fig:performance}
\end{figure*}

TopoStake is not designed to replace the underlying PoS protocol or to increase raw consensus throughput by changing the consensus rule. Its purpose is to add accountable relay incentives while preserving the performance envelope of the base protocol. We therefore compare PoS and TopoStake under increasing transaction loads and network scales.

Fig.~\ref{fig:throughput_rate} and Fig.~\ref{fig:latency_rate} show the results under different input rates. As the input rate increases from 50 to 200 transactions per second, both protocols process more transactions, but TopoStake consistently achieves slightly higher throughput than PoS. At the same time, TopoStake keeps lower p95 confirmation latency across all loads. This indicates that the additional path accounting and reward computation do not introduce prohibitive end-to-end overhead. Instead, the relay incentive helps sustain transaction dissemination under higher load.

Fig.~\ref{fig:throughput_validators} and Fig.~\ref{fig:latency_validators} report scalability with respect to the validator population. Throughput decreases for both protocols as the number of validators grows, reflecting the larger communication and verification workload. However, TopoStake maintains higher throughput than PoS across all tested network sizes. The latency gap is more pronounced: when the validator population increases, the p95 latency of PoS grows rapidly, whereas TopoStake increases more slowly. This suggests that accountable relay incentives can improve propagation efficiency and reduce confirmation delay in larger networks.

Overall, these results show that TopoStake preserves the performance envelope of the base PoS protocol while improving transaction propagation under load and scale. The benefit should be interpreted as an empirical propagation effect rather than a change to the formal liveness threshold of the underlying PoS protocol.

\subsection{Cryptographic Overhead}
\label{subsec:eval_crypto}

\begin{table}[t]
\centering
\caption{BLS microbenchmark under different path lengths. Values are mean runtime in milliseconds.}
\label{tab:crypto_overhead}
\begin{tabular}{c|c|c|c|c}
\hline
Path length & Sign & Verify & Aggregate & Aggregate verify \\
\hline
1  & 0.100 & 0.567  & 1.249  & 0.686 \\
2  & 0.190 & 1.076  & 2.435  & 0.770 \\
4  & 0.385 & 2.195  & 4.820  & 0.874 \\
8  & 0.727 & 4.450  & 9.561  & 0.995 \\
16 & 1.617 & 10.287 & 21.786 & 1.473 \\
\hline
\end{tabular}
\end{table}

TopoStake uses two-sided BLS certificates to make recorded propagation edges publicly verifiable. 
Table~\ref{tab:crypto_overhead} reports the mean cryptographic runtime for paths of 1 to 16 hops. 
The Sign, Verify, and Aggregate columns measure the total per-path cost of generating, checking, and aggregating all sender and receiver edge proofs, while Aggregate verify measures the cost of checking the final aggregate proof included in a block.

The results show that signing, individual verification, and aggregation increase with path length. 
This is expected because each additional hop introduces additional sender and receiver proofs. 
For example, when the path length increases from 1 to 16 hops, individual verification increases from 0.567 ms to 10.287 ms, and aggregation increases from 1.249 ms to 21.786 ms. 
These costs are mainly incurred when constructing and validating path evidence before block inclusion.

In contrast, aggregate verification remains relatively small. 
Even for a 16-hop path, verifying the aggregate proof takes 1.473 ms, which is much lower than individually verifying all edge proofs. 
Therefore, BLS aggregation keeps block-level proof verification lightweight, although the validator sequence and edge-statement reconstruction still grow linearly with the path length.

\subsection{Real-Stake-Bounded Proposer Influence}
\label{subsec:eval_real_stake_bound}

\begin{figure*}[!t]
    \centering
    \subfloat[Score share, $\eta b_{\max}=0.25$\label{fig:score_share_eta025}]{
        \includegraphics[width=0.31\textwidth]{figures/real_stake_bound_score_share_eta0p25.pdf}
    }
    \hfill
    \subfloat[Weight bound, $\eta b_{\max}=0.25$\label{fig:weight_bound_eta025}]{
        \includegraphics[width=0.31\textwidth]{figures/real_stake_bound_weight_bound_eta0p25.pdf}
    }
    \hfill
    \subfloat[Amplification, $\eta b_{\max}=0.25$\label{fig:amp_eta025}]{
        \includegraphics[width=0.31\textwidth]{figures/real_stake_bound_amplification_ratio_eta0p25.pdf}
    }

    \vspace{0.8em}

    \subfloat[Score share, $\eta b_{\max}=0.5$\label{fig:score_share_eta05}]{
        \includegraphics[width=0.31\textwidth]{figures/real_stake_bound_score_share_eta0p5.pdf}
    }
    \hfill
    \subfloat[Weight bound, $\eta b_{\max}=0.5$\label{fig:weight_bound_eta05}]{
        \includegraphics[width=0.31\textwidth]{figures/real_stake_bound_weight_bound_eta0p5.pdf}
    }
    \hfill
    \subfloat[Amplification, $\eta b_{\max}=0.5$\label{fig:amp_eta05}]{
        \includegraphics[width=0.31\textwidth]{figures/real_stake_bound_amplification_ratio_eta0p5.pdf}
    }

    \caption{Real-stake-bounded proposer influence under different adversarial placements and incentive strengths.}
    \label{fig:real_stake_bound}
\end{figure*}

The main security question is whether a low-stake adversary with an advantageous network position can convert propagation centrality into excessive proposer influence.
We evaluate this question by fixing the adversary's real-stake share to 5\%, 10\%, 20\%, and 30\%, and placing adversarial validators either uniformly at random, at high-degree positions, or at high-betweenness positions in the propagation graph.
For each configuration, we measure the adversarial propagation-score share, proposer-weight share, and the theoretical upper bound derived from Theorem~\ref{thm:bounded_influence}.
We run the experiment under two incentive strengths, corresponding to $\eta b_{\max}=0.25$ and $\eta b_{\max}=0.5$; the latter is the largest value allowed by Corollary~\ref{cor:honest_proposer_majority}.

Fig.~\ref{fig:real_stake_bound} shows the results.
The score-share results show that network position does matter: adversaries placed at high-degree or high-betweenness validators obtain substantially larger propagation-score shares than randomly placed adversaries.
This confirms the threat model assumption that propagation contribution is not necessarily proportional to real stake.
However, the proposer-weight results show that this propagation advantage does not translate into unbounded proposer influence.
Across all tested stake fractions, placements, and incentive strengths, the adversarial proposer-weight share remains below the theoretical real-stake-derived bound.

The amplification-ratio results further illustrate the effect of the cap.
For random placement, the ratio $f_{\mathcal{A}}^{W}/f_{\mathcal{A}}^{S}$ stays close to one, while high-degree and high-betweenness placements obtain a moderate amplification due to their larger propagation scores.
Nevertheless, the ratio remains below the protocol cap $1+\eta b_{\max}$ in both parameter settings.
No bound violation is observed in any run.
These results support the main safety claim of TopoStake: propagation scores can reward useful relay positions, but their effect on proposer selection remains bounded by real stake and protocol parameters.

\subsection{Proposer-Side Resilience under Sustained Outages}
\label{subsec:eval_sustained_outage}

\begin{figure*}[!t]
    \centering
    \subfloat[Offline-group proposer-weight response.\label{fig:sustained_outage_weight}]{
        \includegraphics[width=0.47\textwidth]{figures/frozen_v1_sustained_outage/sustained_outage_weight_response.pdf}
    }
    \hfill
    \subfloat[Paired missed-slot-rate improvement.\label{fig:sustained_outage_miss}]{
        \includegraphics[width=0.47\textwidth]{figures/frozen_v1_sustained_outage/sustained_outage_missed_slots.pdf}
    }
    \caption{Proposer-side response to sustained validator outages. Full TopoStake is compared with the paired $\eta=0$ baseline over 20 seeds per condition and 150 outage slots per run; error bars show 95\% mean intervals. Random outages show the dependability benefit after score decay, while high-score outages expose the response delay caused by stale bonuses.}
    \label{fig:sustained_outage}
\end{figure*}

Transaction-scoped relay fees compensate completed propagation work but cannot change future proposer scheduling once a transaction has settled. We therefore isolate the incremental effect of the persistent proposer bonus by comparing Full TopoStake with TopoStake-$\eta0$. Both variants use the same seed-specific outage assignment and randomness schedule. For each seed, a separate 20-epoch calibration fixes the outage assignment; the paired runs then use a 20-epoch warmup followed by a 30-epoch outage (150 slots). Validators controlling a target of 10\%, 25\%, or 40\% of economic stake stop relaying and producing blocks during that outage. We use 20 seeds and consider two assignment rules: \emph{random}, modelling an exogenous outage independent of prior contribution, and \emph{high-score}, a stress case in which validators with the largest pre-outage bonuses fail abruptly. The audit confirms complete paired runs, identical assignments, enforced outages, silent outage relays, a stake-consistent baseline under $\eta=0$, and no proposer-weight-envelope violation.

Fig.~\ref{fig:sustained_outage_weight} reports the outage group's normalized proposer-weight shift relative to TopoStake-$\eta0$. Under random outages, the groups begin near the stake baseline and their steady-state weight shares fall by 0.91, 1.82, and 1.93 percentage points at the 10\%, 25\%, and 40\% targets, respectively. The corresponding missed-slot-rate reductions in Fig.~\ref{fig:sustained_outage_miss} are 2.13, 1.70, and 3.23 percentage points. Thus, once score decay takes effect, the persistent bonus shifts proposer opportunities away from validators that cease producing certified relay activity; a fee-only mechanism cannot provide this future scheduling response.

The high-score stress case exposes the limitation of this signal. At outage onset, stale bonuses raise the failed groups' proposer-weight shares by 1.46--2.45 percentage points relative to the stake-only baseline. After decay, their steady-state shares fall by only 0.11--0.53 percentage points, and the missed-slot intervals overlap zero; at the 40\% target, the mean change is $-1.33$ percentage points. TopoStake therefore does not act as an online failure detector or provide an additional consensus-liveness guarantee beyond the underlying PoS protocol. The result instead demonstrates a bounded, delayed improvement in proposer-side chain-growth resilience under sustained absence of certified relay activity, while finality voting remains stake-only.

\subsection{Path Padding and Transaction Flooding}
\label{subsec:eval_attacks}

\begin{figure*}[!t]
    \centering
    \subfloat[Padding reward\label{fig:path_padding_reward}]{
        \includegraphics[width=0.234\textwidth]{figures/path_padding_coalition_relay_reward.pdf}
    }
    \hfill
    \subfloat[Padding proposer weight\label{fig:path_padding_proposer_weight}]{
        \includegraphics[width=0.234\textwidth]{figures/path_padding_proposal_probability.pdf}
    }
    \hfill
    \subfloat[Flooding net income\label{fig:flooding_net_income}]{
        \includegraphics[width=0.234\textwidth]{figures/transaction_flooding_net_income.pdf}
    }
    \hfill
    \subfloat[Flooding proposer weight\label{fig:flooding_proposer_weight}]{
        \includegraphics[width=0.234\textwidth]{figures/transaction_flooding_proposer_weight.pdf}
    }

    \caption{Robustness of TopoStake under path-padding and transaction-flooding attacks.}
    \label{fig:attack_robustness}
\end{figure*}

We next stress-test TopoStake under path-padding and transaction-flooding attacks.
In the path-padding attack, a malicious proposer and colluding relayers insert additional controlled identities into the recorded propagation path.
We vary the number of padded identities from 0 to 16 and evaluate the coalition's relay reward and proposer-weight share under different target depths.
In the transaction-flooding attack, the adversary generates additional transactions to buy relay credit.
We vary the attack transaction multiplier and measure fee expenditure, recovered relay income, net income, and proposer-weight share.

Fig.~\ref{fig:attack_robustness} reports the results.
Fig.~\ref{fig:path_padding_reward} shows that padding does not increase the coalition's relay reward.
Instead, the reward share decreases rapidly as more controlled identities are inserted, and it approaches zero when the path becomes heavily padded.
This is consistent with the path-budget design: longer paths receive a smaller total budget, so additional identities do not create additional reward capacity.

Fig.~\ref{fig:path_padding_proposer_weight} further shows that padding does not create future proposer advantage.
As the number of padded identities increases, the coalition's proposer-weight share decreases toward its real-stake share.
Thus, path padding is ineffective both for immediate relay reward and for future proposer influence.

Fig.~\ref{fig:flooding_net_income} evaluates transaction flooding.
Although the adversary recovers a small amount of relay income, its fee expenditure grows much faster, and the resulting net income remains negative for all nonzero flooding multipliers.
Fig.~\ref{fig:flooding_proposer_weight} shows that flooding also fails to create unbounded proposer influence: the adversarial proposer-weight share stays close to the real-stake share and remains well below the theoretical bound.
Together, these results support the intended security boundary of TopoStake: syntactically valid colluding paths and self-generated transactions may obtain limited relay credit, but they cannot create positive fee profit or bypass the real-stake-bounded proposer influence.

\subsection{Reward Fairness and Relay Participation}
\label{subsec:eval_fairness}

Finally, we evaluate whether TopoStake rewards active relayers without replacing economic stake as the root security resource. We vary the stake-distribution Gini coefficient and compare PoS, TopoStake-$\eta0$, and TopoStake. We measure reward Gini, HHI, relay income, proposer income, propagation-score distribution, and proposer-weight concentration.

This experiment distinguishes two effects. The relay-fee mechanism should compensate validators that actively forward transactions, including low-stake validators that contribute to dissemination. The proposer-bonus mechanism should provide an additional long-term incentive, but its effect on proposer selection must remain bounded by the protocol parameters. Together, these measurements show whether TopoStake improves relay participation while preserving the real-stake safety boundary.

\section{Conclusion}
The conclusion goes here.


\section*{Acknowledgments}
This should be a simple paragraph before the References to thank those individuals and institutions who have supported your work on this article.




\bibliographystyle{IEEEtran}
\bibliography{reference}




 





\end{document}
